\documentclass[UTF8]{ctexart}
\usepackage{amsmath}
\usepackage{tocloft}
\usepackage[bigfiles]{pdfbase}
\usepackage{amsthm,amssymb}
\usepackage{booktabs}
\usepackage{graphicx}
\usepackage[hidelinks]{hyperref}
\usepackage{geometry}
\usepackage{float}
\usepackage{multirow}
\usepackage{indentfirst}
\usepackage{diagbox}
\usepackage{listings}
\usepackage{xcolor}
\definecolor{codegreen}{rgb}{0,0.6,0}
\definecolor{codegray}{rgb}{0.5,0.5,0.5}
\definecolor{codepurple}{rgb}{0.58,0,0.82}
\definecolor{backcolour}{rgb}{0.95,0.95,0.92}

\lstdefinestyle{mystyle}{
    backgroundcolor=\color{backcolour},   
    commentstyle=\color{codegreen},
    keywordstyle=\color{magenta},
    numberstyle=\tiny\color{codegray},
    stringstyle=\color{codepurple},
    basicstyle=\ttfamily\footnotesize,
    breakatwhitespace=false,         
    breaklines=true,                 
    captionpos=b,                    
    keepspaces=true,                 
    numbers=left,                    
    numbersep=5pt,                  
    showspaces=false,                
    showstringspaces=false,
    showtabs=false,                  
    tabsize=2
}
\lstset{style=mystyle}
\newcommand\question[2]{\noindent\fbox{\parbox[t]{\linewidth}{\hspace{0pt}{\textbf{#1\quad}#2}}}}


\usepackage{graphicx}
\usepackage{setspace}
\renewcommand{\cftsecleader}{\cftdotfill{\cftdotsep}}
\geometry{a4paper,left=3cm,right=3cm,top=2cm,bottom=2cm} 


\CTEXsetup[format={\Large \bfseries}]{section}
\title{\textbf{神经网络加速器计算核心设计}}
\author{陈天乐\ \ 无25\ \ [student ID removed] }
\date{\today}
\begin{document}
\maketitle
\tableofcontents
\newpage
\section{实验目的}
参考“加速器”、“映射”部分的课程内容，使用 Verilog/SystemVerilog语言设计实现 CNN 加速器的核心加速模块——计算阵列。
\section{基本任务}

\subsection{乘累加树型阵列设计与实现}

本实验首先实现了一个 $1\times N$ 乘累加树阵列，对应代码文件为
\verb|mac_tree.sv|，测试文件为 \verb|tb_mac_tree.sv|。
阵列主要参数为 \verb|MAC_NUM| 和 \verb|BRANCH_NUM|。其中
\verb|MAC_NUM| 表示并行输出lane的数量；\verb|BRANCH_NUM| 表示每棵加法树
在一个周期内接收的输入激活值和权重对数。取$N=10$，因此每个输出块一共需要15个计算周期完成累加。

该设计的前仿结果完全正确，总耗时为25645ns，仿真结果如下：

\begin{figure}[H]
    \centering
    \includegraphics[width=1\textwidth]{1.png}
    \caption{1*10乘累加树计算阵列仿真波形}
\end{figure}

\begin{lstlisting}[language={},caption={乘累加树阵列仿真检查结果}]
mac_tree match_count=1600/1600, wrong_num=0
PASS: mac_tree all 1600 outputs match reference.
\end{lstlisting}

乘累加树型阵列的资源利用如下：
\begin{figure}[H]
    \centering
    \includegraphics[width=0.6\textwidth]{3.png}
    \caption{1*10乘累加树计算阵列资源利用报告}
\end{figure}

乘累加树型阵列的时序报告如下：
\begin{figure}[H]
    \centering
    \includegraphics[width=0.6\textwidth]{2.png}
    \caption{1*10乘累加树计算阵列时序报告}
\end{figure}

通过时序报告可以看到，时序约束符合要求，并且在100MHz的时钟下，WNS为8.477ns，说明可以用更快的时钟运行，并且提升空间很高。
\begin{table}[H]
    \centering
    \caption{乘累加树阵列在 100MHz 下的设计参数}
    \begin{tabular}{cc}
        \toprule
        算力 & 10*10*100M=10Gops \\
        计算量/算力 & 100*150*16/10Gops=24us \\
        \midrule
        最高时钟频率 & $\frac{1}{10ns-8.477ns}=656.60MHz$ \\
        \midrule
        阵列利用率 & 24/25.645 = 93.59\% \\        
        \midrule
        IA带宽需求 & $8\mathrm{b}\times10\times10\times100\mathrm{MHz}=80\mathrm{Gbps}$ \\
        W带宽需求 & $8\mathrm{b}\times10\times100\mathrm{MHz}=8\mathrm{Gbps}$ \\
        OA带宽需求 & $8\mathrm{b}\times10\times100\mathrm{MHz}=8\mathrm{Gbps}$ \\
        \midrule
        IA 片上 buffer 容量 & $8\mathrm{b}\times10\times10=800\mathrm{bit}$ \\
        W 片上 buffer 容量 & $8\mathrm{b}\times10=80\mathrm{bit}$ \\
        OA 片上 buffer 容量 & $8\mathrm{b}\times10=80\mathrm{bit}$ \\
        \midrule
        IA 片外访问总 bit 数 & $16\times100\times150\times8\mathrm{bit}$ \\
        OA 片外写回总 bit 数 & $16\times10\times150\times8\mathrm{bit}$ \\
        W 片外访问总 bit 数 & $16\times100\times8\mathrm{bit}$ \\
        \bottomrule
    \end{tabular}
\end{table}

\subsection{脉动阵列型设计与实现}

第二个实验实现了 $1\times N$ 脉动阵列，对应代码文件为
\verb|systolic_array.sv|，测试文件为 \verb|tb_systolic_array.sv|。
我设计的脉动阵列的\verb|CELL_H=10|，\verb|CELL_W=1|。阵列中 10 个 PE 纵向连接，采用weight-stationary策略，权重在开始时加载到 PE 内部并保持不变，激活值按照斜向时序输入，部分和沿阵列方向
逐级向下传递。这样右侧或下游 PE 的乘加结果会比上游晚一个周期出现，符合脉动阵列
数据逐拍传播的特征。需要指出的是，第 0 个 PE 接收外部部分和
    \verb|PE_partial_sum|，后续 PE 接收上一级 PE 的结果，形成一条 10 级流水。
仿真检查结果如下：



该设计的前仿结果完全正确，总耗时为266445ns，主要开销来自15 个 tile需要依次流过整条 10 级脉动链。仿真结果如下：

\begin{figure}[H]
    \centering
    \includegraphics[width=1\textwidth]{4.png}
    \caption{1*10脉动阵列仿真波形}
\end{figure}

\begin{lstlisting}[language={},caption={$1\times 10$ 脉动阵列仿真检查结果}]
systolic_array match_count=1600/1600, wrong_num=0
PASS: systolic_array all 1600 outputs match reference.
\end{lstlisting}

1*10脉动阵列的资源利用如下：
\begin{figure}[H]
    \centering
    \includegraphics[width=0.6\textwidth]{6.png}
    \caption{1*10脉动阵列资源利用报告}
\end{figure}

1*10脉动阵列的时序报告如下：
\begin{figure}[H]
    \centering
    \includegraphics[width=0.6\textwidth]{5.png}
    \caption{1*10脉动阵列时序报告}
\end{figure}

通过时序报告可以看到，时序约束符合要求，并且在100MHz的时钟下，WNS为5.727ns，说明可以用更快的时钟运行，并且提升空间很高。
\begin{table}[H]
    \centering
    \caption{1*10脉动阵列在 100MHz 下的设计参数}
    \begin{tabular}{cc}
        \toprule
        算力 & 1*10*100M=1Gops \\
        计算量/算力 & 100*150*16/1Gops=240us \\
        \midrule
        最高时钟频率 & $\frac{1}{10ns-5.727ns}=234.02MHz$ \\
        \midrule
        阵列利用率 & 240/266.445 = 90.07\% \\        
        \midrule
        IA带宽需求 & $8\mathrm{b}\times1\times10\times100\mathrm{MHz}=8\mathrm{Gbps}$ \\
        W带宽需求 & $8\mathrm{b}\times10\times100\mathrm{MHz}=8\mathrm{Gbps}$ \\
        OA带宽需求 & $32\mathrm{b}\times100\mathrm{MHz}=3.2\mathrm{Gbps}$ \\
        \midrule
        W 片上 buffer 容量 & $8\mathrm{b}\times10=80\mathrm{bit}$ \\
        OA 片上 buffer 容量 & $32\mathrm{bit}$ \\
        \midrule
        IA 片外访问总 bit 数 & $16\times100\times150\times8\mathrm{bit}$ \\
        OA 片外写回总 bit 数 & $16\times10\times150\times8\mathrm{bit}$ \\
        W 片外访问总 bit 数 & $16\times100\times8\mathrm{bit}$ \\
        \bottomrule
    \end{tabular}
\end{table}


在 100MHz 时钟下，该阵列每拍完成 10 个乘加，理论峰值算力为
$10\times100\mathrm{MHz}=1\mathrm{Gops}$。由于只有一个输出列，输出写回峰值带宽为
$8\mathrm{b}\times100\mathrm{MHz}=800\mathrm{Mbps}$，低于参考 $1\times10$ 独立 MAC
阵列的 8Gbps 输出压力。
\subsection{拓展脉动阵列为 $M\times N$ 设计与实现}

本实验将上述1*N脉动阵列扩展为 $M\times N$，对应代码文件为
\verb|multi_systolic_array.sv|，测试文件为
\verb|tb_multi_systolic_array.sv|。本次参数取 \verb|CELL_H=10|、
\verb|CELL_W=4|，一次并行计算 4个输出通道，每个输出通道对应一列 partial sum
流水链。由于本次输出通道数为 16，刚好可以被 4 整除，因此 testbench 每次处理 4 个
输出通道，循环 4 次即可覆盖全部输出通道。

拓展的关键是将 $1\times N$ 脉动阵列中的部分和寄存器扩展为二维：
\verb|psum_reg[CELL_H][CELL_W]|。同一个激活值同时广播给 4 个输出列，而每一列使用
各自的权重和 partial sum，因此可以复用 IA 读取并提升输出通道方向的并行度。对于
\verb|CELL_W=4| 的情况，每个tile 加载 $10\times4$ 个权重，随后流过 100 行输入，
同时更新 4 个输出通道的 100 个 partial sum。

该设计的前仿结果完全正确，总耗时为66645ns，与上一设计时间开销缩小为1/4，符合预期。仿真结果如下：


\begin{figure}[H]
    \centering
    \includegraphics[width=1\textwidth]{4.png}
    \caption{4*10脉动阵列仿真波形}
\end{figure}

\begin{lstlisting}[language={},caption={$4\times 10$ 脉动阵列仿真检查结果}]
multi_systolic_array match_count=1600/1600, wrong_num=0
PASS: multi_systolic_array all 1600 outputs match reference.
\end{lstlisting}




1*10脉动阵列的资源利用如下：
\begin{figure}[H]
    \centering
    \includegraphics[width=0.6\textwidth]{8.png}
    \caption{4*10脉动阵列资源利用报告}
\end{figure}

1*10脉动阵列的时序报告如下：
\begin{figure}[H]
    \centering
    \includegraphics[width=0.6\textwidth]{9.png}
    \caption{4*10脉动阵列时序报告}
\end{figure}

通过时序报告可以看到，时序约束符合要求，并且在100MHz的时钟下，WNS为5.727ns，说明可以用更快的时钟运行，并且提升空间很高。
\begin{table}[H]
    \centering
    \caption{4*10脉动阵列在 100MHz 下的设计参数}
    \begin{tabular}{cc}
        \toprule
        算力 & 4*10*100M=4Gops \\
        计算量/算力 & 100*150*16/4Gops=60us \\
        \midrule
        最高时钟频率 & $\frac{1}{10ns-5.727ns}=234.02MHz$ \\
        \midrule
        阵列利用率 & 60/66.645 = 90.03\% \\        
        \midrule
        IA带宽需求 & $8\mathrm{b}\times1\times10\times100\mathrm{MHz}=8\mathrm{Gbps}$ \\
        W带宽需求 & $8\mathrm{b}\times10\times100\mathrm{MHz}\times4=32\mathrm{Gbps}$ \\
        OA带宽需求 & $32\mathrm{b}\times100\mathrm{MHz}\times4=12.8\mathrm{Gbps}$ \\
        \midrule
        W 片上 buffer 容量 & $8\mathrm{b}\times10\times4=320\mathrm{bit}$ \\
        OA 片上 buffer 容量 & $32\mathrm{b}*4=128\mathrm{bit}$ \\
        \midrule
        IA 片外访问总 bit 数 & $16\times100\times150\times8\mathrm{bit}$ \\
        OA 片外写回总 bit 数 & $16\times10\times150\times8\mathrm{bit}$ \\
        W 片外访问总 bit 数 & $16\times100\times8\mathrm{bit}$ \\
        \bottomrule
    \end{tabular}
\end{table}


在 100MHz 时钟下，该阵列每拍完成 $10\times4=40$ 个乘加。相对于 $1\times N$ 脉动阵列，IA 带宽基本
不变，但同一份输入激活值被 4 个输出通道共享。


从结果看，乘累加树适合在 $K$ 维上展开并行，单次输出块的计算周期较少，但 IA 输入峰值
带宽较高；$1\times N$ 脉动阵列输出带宽压力较小，但输出通道并行度不足；$M\times N$
脉动阵列在复用同一份 IA 的同时并行计算多个输出通道，是本实验中综合延时和带宽复用
更均衡的方案。

\section{思考题}

\question{1.}{用非 GEMM 的方式算卷积（比如 Winograd 卷积，或者直接算卷积），该如何控制数据向阵列的传输？}

对于 Winograd 卷积，输入特征图和卷积核都不能直接按普通 GEMM 的 $K$ 维输入阵列，而是
需要先按 tile 做变换。例如常见的 $F(2\times2,3\times3)$ 会把一个输入小块变换到
Winograd 域，再与变换后的权重做逐点乘，最后做反变换得到输出。因此数据控制上应当
增加一个 tile buffer：先从片外读入覆盖当前输出 tile 所需的输入区域，在片上完成
Winograd 输入变换，再将变换后的若干通道数据送入乘加阵列。权重也可以预先变换后缓存在
片上，减少运行时重复变换。

对于直接卷积，可以不做 img2col 展开，而是在片上维护 line buffer 和滑动窗口 buffer。
当卷积窗口向右移动时，大部分输入像素可以复用，只需要读入新进入窗口的一列；当窗口
换行时，line buffer 更新一行数据。这样送入阵列的不是完整展开矩阵，而是窗口内的
$C\times R\times S$ 个激活值和对应权重。控制逻辑需要产生窗口坐标、通道坐标和 padding
判断，并将同一个输入像素广播到多个输出通道的 PE 中，提高输入复用率。\\

\question{2.}{在你之前设计的基础上，如果有更大的片上 buffer（或者有一个新的 shared buffer），
你会怎么重新进行 mapping 的设计？会有什么样的好处？（提示：40nm 工艺下，外部访存
的功耗是内部访存的 3$\sim $6 倍以上）可以进行一个 case study。}

如果有更大的片上 buffer，我会优先改变权重和 partial sum 的 mapping，使其尽量留在片上。
本次任务的权重矩阵大小为 $150\times16\times8\mathrm{bit}=19200\mathrm{bit}$，约为
2.4KB，实际上很适合完整缓存在 shared buffer 中。若将全部权重一次性读入片上，则乘累加树
mapping 中权重不必对每个 10 行输出块重复从片外读取，W 片外访问量可以从
$16\times10\times150\times8\mathrm{bit}$ 降到
$16\times150\times8\mathrm{bit}$，约减少 10 倍。

对于 $M\times N$ 脉动阵列，还可以在 shared buffer 中保存若干输出行、若干输出通道的
32 位 partial sum。这样每个 $K$ tile 结束后不必把 partial sum 写回片外，而是在片上等待
下一个 $K$ tile 继续累加，最终只写回 8 位量化后的 OA。由于片外访存功耗显著高于片内
访存，这种 weight/partial-sum stationary 的 mapping 能同时降低能耗和延时。若 buffer
足够大，还可以采用 double buffer，在阵列计算当前 tile 时预取下一 tile 的 IA 和 W，隐藏
数据搬运延迟。\\

\question{3.}{在你之前设计的基础上，如果实际不能满足你的带宽需求，比如参考代码的 OA 带宽需求
是 8000Mbps，但实际只能做到 800Mbps，那该如何重新进行设计？可以进行一个 case study。}

如果 OA 带宽只有 800Mbps，本质上等价于 100MHz 下每周期只能写回 1 个 8 位输出。对于
原参考代码或本实验的乘累加树阵列，一次可能产生 10 个 8 位输出，峰值 OA 带宽为
8Gbps，不能直接写回片外。解决方法是增加片上 output buffer，将阵列一次产生的多个输出
先写入片上 FIFO，再用 1 个 8 位端口串行写回片外。

以本实验的乘累加树为例，每个输出块得到 10 个输出。如果外部 OA 端口只有 800Mbps，则
需要 10 个周期才能写完这 10 个结果。为了避免阵列完全停顿，可以使用双 output buffer：
buffer A 串行写回上一组输出时，阵列继续计算下一组输出并写入 buffer B。当计算时间短于
写回时间时，最终性能仍会被 OA 带宽限制；此时应减少输出并行度，或者改用 $1\times N$
脉动阵列这类每周期只产生 1 个输出的结构。本实验的 $1\times N$ 脉动阵列输出写回峰值
正好是 $8\mathrm{b}\times100\mathrm{MHz}=800\mathrm{Mbps}$，因此在该带宽约束下更容易
稳定运行，只是算力利用率较低。\\

\question{4.}{在你之前设计的基础上，如果阵列的尺寸并不是新任务中输入矩阵尺寸的因数，例如在
参考代码中，阵列尺寸是 1*10，来了个新任务输入激活值矩阵尺寸是 102*128，权重矩阵尺
寸是 128*18。那么你该如何进行 mapping？或者如何进行更好的阵列设计？可以进行一个case study。}

若仍使用 $1\times10$ 阵列，可以采用分块加补零的方式完成 mapping。输入矩阵为
$102\times128$，权重矩阵为 $128\times18$，输出为 $102\times18$。行方向按 10 行切分时，
前 10 个 tile 覆盖 100 行，最后一个 tile 只有 2 行有效，需要补 8 行零；输出通道方向按
1 个通道循环 18 次；$K=128$ 若按 10 切分，则需要 13 个 $K$ tile，最后一个 tile 只有
8 个有效元素，剩余 2 个补零。这样可以保证结果正确，但最后的行 tile 和最后的 $K$ tile
都会降低阵列利用率。

更好的阵列设计应尽量让阵列参数贴合任务尺寸。例如可以选择输出通道并行度 $M=6$，因为
18 可以被 6 整除；$K$ 方向可以选择 8 或 16，因为 128 可以被 8 和 16 整除。对于脉动阵列，
一个 $8\times6$ 或 $16\times6$ 的阵列都比 $10\times1$ 更适合该任务：权重 tile 不需要
补零，输出通道也没有尾块浪费。若还希望行方向也整除，可以让乘累加树的 row lane 取
6 或 17 等能更好覆盖 102 的参数，或者采用可屏蔽 lane 的设计，在最后一个 tile 只使能
有效输出行。\\

\question{5.}{在你之前设计的基础上，试着去搜索一下更好的 mapping 方案。（延时更低、带宽需求更
低、访问 bit 数更低等）}

更好的 mapping 通常围绕数据复用展开。本实验中比较直接的优化是采用 weight stationary
或 output stationary 的混合策略：权重矩阵较小，可以完整缓存在片上；partial sum 使用
32 位片上 buffer 保存，直到所有 $K$ tile 累加结束后再量化写回。这样可以减少权重和
中间结果的片外访问。对于卷积本身，如果不强制 GEMM 展开，还可以采用 row stationary
mapping：输入特征图的相邻卷积窗口共享大量像素，使用 line buffer 保存多行输入，使每个
输入像素在多个窗口和多个输出通道之间复用，从而降低 IA 访问 bit 数。

对于本实验的 $M\times N$ 脉动阵列，当前已经体现了一部分输入复用：同一个
\verb|PE_act_in[i]| 同时广播给 4 个输出通道列。若继续增大 \verb|CELL_W|，例如从 4 增大到
8 或 16，可以进一步提高权重列方向并行度，并减少重复读取 IA 的次数。但 \verb|CELL_W|
增大后，权重带宽和 partial sum 寄存器数量也会增加，因此需要结合片上 buffer 容量和外部
带宽共同选择。\\

\question{6.}{试着优化你的设计，使其能够跑到更高的时钟/更短的任务时间。}

要提高时钟频率，乘累加树设计的主要瓶颈是组合乘法和多级加法树路径较长。可以在加法树
各层之间加入流水寄存器，或者使用 FPGA 的 DSP 乘法器内部寄存器，把一拍内的长组合路径
拆成多拍流水。这样会增加输出延迟，但可以显著提高最高时钟频率；testbench 只需要相应
增加结果采样延迟即可。

要缩短任务时间，可以从三个方向优化。第一，增大并行度，例如将 $M\times N$ 脉动阵列的
\verb|CELL_W| 从 4 提高到 8 或 16，使一次处理更多输出通道。第二，减少控制气泡，本实验中
每个 $K$ tile 之间有加载权重、清空流水等控制周期，可以用双缓冲权重寄存器和 valid 信号
让下一组权重预加载，从而让计算流更连续。第三，优化片上 buffer，使权重和 partial sum
尽量不访问片外，并通过 double buffer 重叠数据搬运和计算。综合来看，若资源允许，采用
更宽的 $M\times N$ 脉动阵列并配合片上 weight buffer 和 partial sum buffer，是本任务中
获得更短运行时间的主要方向。

\end{document}
